% Lösungen Einheit 1 von 4 – Geradengleichung aufstellen und lesen (zusammenbau v0.1)
\einheitenkopf{Einheit 1 von 4 · Geradengleichung aufstellen und lesen}

\erg{10}{a) Strecke – $r = 0$ ist der Anfang des Förderbands am Boden, $r = 1$ sein oberes Ende.}
\erg{11}{a) $(3 | 0 | -2)$ – der Vektor, der mit dem Parameter $t$ multipliziert wird \quad b) $g\colon \vec x = (2 | 1 | 3) + r \cdot (2 | 3 | -1)$ \quad c) $\vec x = (2 | -4 | 1) + \lambda \cdot (-8 | 24 | 12)$ mit $0 \le \lambda \le 1$: die Liftstrecke von der Talstation ($\lambda = 0$) bis zur Bergstation ($\lambda = 1$), nicht die ganze Gerade. \quad d) $\overrightarrow{AB} = (0 | 6 | 0)$: erste und dritte Koordinate stimmen überein, nur die zweite ändert sich – die Gerade verläuft parallel zur $x_2$-Achse (sie liegt nicht auf ihr). \quad e) zum Beispiel $p\colon \vec x = (4 | 1 | 2) + r \cdot (1 | 3 | -2)$ – der Stützpunkt liegt nicht auf $g$, weil $(1 | 0 | 0)$ kein Vielfaches von $(1 | 3 | -2)$ ist; $q\colon \vec x = (3 | 1 | 2) + r \cdot (2 | 0 | 1)$ – Stützpunkt von $g$ und $2 + 0 - 2 = 0$. \quad f) $T = P + \tfrac{1}{4} \cdot \overrightarrow{PQ} = (4 | 0 | 2)$, also $\vec x = (4 | 0 | 2) + r \cdot (1 | 4 | 2)$ \quad g) Richtungsvektor $1 \cdot (1 | 0 | 2) + 2 \cdot (0 | 3 | -1) = (1 | 6 | 0)$ mit dritter Komponente null, also $g\colon \vec x = (3 | 2 | 2) + \lambda \cdot (1 | 6 | 0)$}
\erg{12}{a) $(3 | 5 | 1)$ liegt auf $g$ (für $r = 1$), also ist $p = g$ und nicht echt parallel. Richtig: ein Stützpunkt außerhalb von $g$, etwa $p\colon \vec x = (1 | 4 | 0) + r \cdot (2 | 1 | 3)$. \quad b) Alle Punkte der Geraden haben dieselbe dritte Koordinate wie der Stützpunkt: ist sie nicht null, trifft die Gerade die Ebene nie; ist sie null, liegt die Gerade ganz in der Ebene. \quad c) $\vec x = (0 | 0 | 20) + t \cdot (4 | 3 | 1)$ mit $0 \le t \le 30$; nach 10 s: $(40 | 30 | 30)$ \quad d) Stützpunkt $(2 | 0 | 3)$; für $t = 1$: $(2 | 2 | 2)$; für $t = 2$: $(2 | 4 | 1)$}

12d)

\begin{ksys3} \rgerade{2,0,3}{0,2,-1}{g} \rvektorab{2,0,3}{2,2,2}{\vec u} \rpunkt{2}{0}{3}{P_0} \rpunkt{2}{2}{2}{P_1} \rpunkt{2}{4}{1}{P_2} \end{ksys3}

